% =============================================================================
% Chapter 5 --- Verification Methodology
% =============================================================================

\chapter{Verification Methodology}
\label{ch:verification}

This chapter addresses two questions that a senior hardware reviewer
would reasonably ask of a project of this type: \emph{what was the
verification trying to prove}, and \emph{how was it possible to prove it
without access to the physical analog ANN core?} The answers justify a
third claim, made forward: that the present functional simulation
delivers enough confidence to hand the RTL to the analog phase of the
project. Mechanical testbench details --- file structure, helper
classes, console-log formatting --- are documented in the repository.

\section{Verification Objectives}
\label{sec:verif-goals}

The controller has no datapath arithmetic of its own. What it
\emph{does} is translate host transactions into a stream of analog
operations on a Y-Flash array and react correctly to the analog-side
observables that come back. Verification therefore aims to establish
three classes of claim:

\begin{enumerate}
    \item \textbf{Static transaction-word contract.} For every legal
          host transaction, the controller produces a well-formed core
          word and the correct pulse-mode encoding: address fields
          bit-consistent with the host-supplied one-hot tail, payload
          forwarded faithfully, and the pulse-mode bus carrying the
          encoding appropriate to the active operation.
    \item \textbf{Dynamic protocol compliance.} The FSM transitions
          through its states in the prescribed order, with the
          prescribed handshakes, without stalling. Programming advances
          \texttt{SELECT}~$\rightarrow$~\texttt{WRITE}~$\rightarrow$
          \texttt{WAIT\_ACK}~$\rightarrow$~\texttt{COMPLETE}; verify
          follows programming; erase is entered from the correct path
          (host or verify) and exits to the correct main state;
          inference collects the right number of pixels before
          compute. Every protocol error here would manifest only over
          multiple cycles in interaction across sub-FSMs.
    \item \textbf{Closed-loop recovery correctness.} The central
          architectural claim. Match $\rightarrow$ exit. Under
          $\rightarrow$ retry within budget, fall to erase otherwise.
          Over $\rightarrow$ erase, then retry within budget or
          terminate. Every branch of the recovery diagram in
          \Cref{sec:rtl-controller-recovery} is a verification
          obligation.
\end{enumerate}

\section{The Behavioral Mock: Standing in for the Analog ANN}
\label{sec:verif-mock-strategy}

The analog Y-Flash array is not part of this project, so verification
substitutes it with a behavioral model that exposes the same digital
boundary signals --- the inbound core word and pulse-mode bus, the
outbound readback line and \emph{op\_done} handshake. The mock captures
only the digital contract at the boundary, not the analog physics; that
is the right choice because the controller's correctness depends on those
four signals and on nothing else. The integration is shown in
\Cref{fig:testbench-stack}.

\begin{figure}[h]
    \centering
    \includegraphics[width=0.9\textwidth]{example-image.png}
    \caption{Verification environment. The three RTL modules are
        instantiated in their real host-to-core order; the analog ANN
        array is replaced by a behavioral mock with a shadow weight
        matrix, a combinational readback generator driven by the
        outbound core word, and a clocked \emph{op\_done} generator that
        follows the controller's own sub-FSM. Fault injection at the
        readback line during the verify window forces the
        under- and over-programmed branches of the recovery FSM.}
    \label{fig:testbench-stack}
\end{figure}

The mock has three pieces. The \textbf{shadow weight matrix} is a
four-dimensional array indexed by the same
(\textit{PE}, \textit{SA}, row, column) coordinates as the address tail.
On every cycle in which the controller is in the active programming
phase, the mock decodes the outbound core word and writes the quantized
weight nibble carried in the payload byte into the corresponding shadow
cell. A subsequent read returns that cell's value. The
\textbf{behavioral readback generator} continuously decodes the current
outbound core word and selects the shadow cell whose value should appear
on the readback line; with this default behavior alone, a well-behaved
program always satisfies the verify comparator on the first attempt, so
the recovery branches are never visited --- which motivates the fault
injection of \Cref{sec:verif-injection}. The \textbf{\emph{op\_done}
generator} is a small clocked process that asserts the handshake
whenever the controller is in a sub-state structurally waiting for it
(\texttt{PROG\_SELECT}, \texttt{PROG\_WAIT\_ACK}, \texttt{ERASE\_SELECT},
\texttt{ERASE\_WAIT\_ACK}), with a small delay counter on pulse-driven
phases to prevent the bench from stalling. This is testbench-only
timing: it is not modeling a specific analog latency, only enough to let
the FSM advance.

\paragraph{Why the mock is sufficient for confidence.}
Three observations justify the choice. \emph{The controller only sees
the digital contract}: every other physical attribute of the array ---
voltage levels, leakage, parasitic settling --- is invisible to the FSM,
so any controller bug that depends only on digital observables will be
reproduced in simulation. \emph{The mock matches the addressing
convention exactly}: the shadow matrix is indexed by the same
decomposition the RTL itself uses, eliminating the largest single source
of mock-RTL drift. \emph{The mock is composable}: swapping it for a
higher-fidelity model (a Verilog-A cell model or an analog co-simulation
harness) requires no change to either the controller or the testbench
skeleton. The directed regressions written here will run against any
model that conforms to the same digital boundary.

\section{Forcing the Recovery Branches: Fault Injection}
\label{sec:verif-injection}

The most subtle architectural element of the verification environment is
the fault-injection mechanism. Its purpose is to address the verification
gap that the default mock leaves open: a well-behaved mock always returns
the freshly programmed value, so the under- and over-programmed branches
of the verify sub-FSM would never be visited and the entire closed-loop
policy --- the central architectural claim of the project --- would
remain unverified.

The mechanism overrides the readback value during a carefully gated
window inside the verify phase. The override fires only when three
conditions hold: the current transaction is one of a pre-declared set of
injection indices, the controller is in the early cycles of the verify
read window, and the pulse-mode bus carries a verify-related encoding.
Outside this window the readback line carries the unperturbed shadow
value. Two complementary injection scenarios are applied to disjoint
index sets:

\begin{itemize}
    \item \textbf{Under-programmed injection} drives the readback one
          step \emph{below} the expected weight. The controller sees a
          short-of-target condition, increments the per-cell retry
          counter, and re-enters \texttt{PROGRAM} with the short
          single-cycle retry pulse train. The branches exercised are:
          \texttt{VERIFY\_CHECK}~under, retry-counter increment,
          \texttt{PROG} retry, eventual successful verify (after the
          injection window expires). If the injection persists past the
          retry budget, the FSM falls into the erase branch.
    \item \textbf{Over-programmed injection} drives the readback one
          step \emph{above} the expected weight. The controller takes
          the over-programmed branch, enters \texttt{ERASE} with the
          verify-initiated flag, completes the erase sub-FSM, and then
          either returns to \texttt{PROGRAM} (if the global retry
          counter is still in budget) or terminates. The branches
          exercised are: \texttt{VERIFY\_CHECK}~over,
          verify-initiated erase, full erase sub-FSM,
          \texttt{ERASE\_COMPLETE}~$\rightarrow$~\texttt{PROGRAM},
          and a re-run of the program--verify sequence on the same
          cell.
\end{itemize}

\paragraph{Coverage discipline.}
Each program--verify regression run is constrained to fail unless both
the under- and over-programmed paths have been observed at least once.
This guards against accidental regressions in the injection logic or the
gating window being silently absorbed by a still-passing scoreboard.

\section{End-of-Run Checks}
\label{sec:verif-loop}

Three classes of check close the verification loop. \emph{Static-contract
checks} compare observed outputs against expected values computed by the
same helper functions the RTL uses (so the comparison is a true contract
check rather than a re-implementation). The \emph{shadow-matrix
scoreboard} compares post-run shadow contents against the expected
weight vector; because every successful verify writes the expected value
into the shadow, a clean scoreboard confirms that the recovery policy
converged the matrix to the intended state for every cell. The
\emph{branch-exercise check} asserts that each recovery branch has been
entered at least once, preventing the regression from silently
degenerating into a happy-path check.

\section{Confidence Argument for Analog Integration}
\label{sec:verif-confidence-future}

No claim is made here about analog correctness. What \emph{is} claimed
is that the controller's digital behavior --- every state transition,
every pulse train, every recovery branch, every transaction-word
format --- has been exercised under the same digital boundary contract
that the analog core will eventually expose. The controller is verified
\emph{in situ} (together with the parallel interface and the input
buffer, in real host-to-core order), \emph{by construction} (fault
injection forces every recovery branch to be visited), and
\emph{composably} (a higher-fidelity model can replace the mock without
disturbing the surrounding regressions).

What remains for the analog phase is therefore confined to the analog
domain: device variability, pulse-amplitude sensitivity, IR drop on the
array, and timing of the \emph{op\_done} handshake relative to actual
cell settling. None of those concerns invalidate the digital contract;
they refine it. The functional regressions established here become the
baseline that the analog phase reuses, and any divergence observed in
mixed-signal simulation is then attributable to the new model rather
than to a previously unverified piece of the FSM.
